% !TeX root = part7.tex
\documentclass[../final.tex]{subfiles}

\begin{document}

% Лекция 13.10.2026.
% Источник: Desktop/4.jpg (с отметки 13.10), 5.jpg, 6.jpg, 7.jpg, 8.jpg.
\ifSubfilesClassLoaded{\setcounter{chapter}{5}\setcounter{section}{8}}{}
\captionsetup{font=footnotesize,labelfont=bf,skip=5pt}

\section{Устойчивость явной и неявной схем Эйлера}

\subsection{Тестовое уравнение}

Рассмотрим одну затухающую моду линейной системы:
\begin{equation*}
    u'=-\lambda u,\qquad u(0)=u_0,\qquad \lambda>0.
\end{equation*}
Её точное решение
\begin{equation*}
    u(x)=u_0e^{-\lambda x}
\end{equation*}
стремится к нулю. Возмущение начальных данных также убывает как
$e^{-\lambda x}$. Численный метод должен воспроизводить это затухание.

Для схемы, которая на этом уравнении имеет вид
\begin{equation*}
    y_{i+1}=q(\lambda h)y_i,
\end{equation*}
величина $q$ называется \rconcept{множителем усиления}.
Если $\widetilde y_i=y_i+\delta_i$, то
\begin{equation*}
    \widetilde y_{i+1}=q(y_i+\delta_i),\qquad
    \delta_{i+1}=q\delta_i,\qquad
    \delta_i=q^i\delta_0.
\end{equation*}
Поэтому $|q|\le1$ означает, что начальное возмущение не возрастает,
а при $|q|<1$ оно затухает.

\subsection{Явная схема: ограничение на шаг}

Для явного метода Эйлера
\begin{equation*}
    \frac{y_{i+1}-y_i}{h}=-\lambda y_i
    \quad\Longrightarrow\quad
    y_{i+1}=(1-\lambda h)y_i.
\end{equation*}
Следовательно,
\begin{equation*}
    q_{\text{явн}}=1-\lambda h,\qquad
    \delta_{i+1}=(1-\lambda h)\delta_i.
\end{equation*}
Чтобы ошибка не росла, необходимо
\begin{equation*}
    |1-\lambda h|\le1
    \quad\Longleftrightarrow\quad
    0\le\lambda h\le2.
\end{equation*}
При положительном шаге получаем \rresult{условие устойчивости}
\begin{equation*}
    0<h\le\frac2\lambda.
\end{equation*}
При $0<\lambda h<1$ приближения убывают без смены знака,
при $\lambda h=1$ следующее значение равно нулю, а при
$1<\lambda h<2$ знак чередуется, но модуль убывает.
На границе $\lambda h=2$ множитель равен $-1$: возмущение
не затухает. При $\lambda h>2$ схема неустойчива.

\remark
\rnote{Устойчивость не гарантирует точности}. Даже при $|q|<1$
численное решение может заметно отличаться от $e^{-\lambda x}$.
В частности, чередование знака отсутствует у точного решения
с положительным $u_0$.

\subsection{Условная и безусловная устойчивость}

Схема называется \rconcept{условно устойчивой}, если условие
устойчивости выполняется лишь при ограничении на шаг. Явная схема
Эйлера на рассматриваемом уравнении требует $h\le2/\lambda$.

Если схема устойчива при любом $h>0$ для данного класса задач,
говорят о \rconcept{безусловной устойчивости}. В записи лекции
для этого также используется выражение «абсолютная устойчивость
при любом шаге».

\subsection{Неявная схема Эйлера}

В неявной схеме правая часть вычисляется на новом слое:
\begin{equation*}
    y_{i+1}=y_i+h f(x_{i+1},y_{i+1}).
\end{equation*}
Для тестового уравнения получаем
\begin{align*}
    y_{i+1}&=y_i-\lambda h y_{i+1},\\
    (1+\lambda h)y_{i+1}&=y_i,\\
    y_{i+1}&=\frac{y_i}{1+\lambda h}.
\end{align*}
Возмущение изменяется по закону
\begin{equation*}
    \delta_{i+1}=\frac{\delta_i}{1+\lambda h}.
\end{equation*}
При $\lambda>0$ и любом $h>0$
\begin{equation*}
    0<q_{\text{неявн}}=\frac1{1+\lambda h}<1.
\end{equation*}
Следовательно, неявный метод Эйлера
\rresult{безусловно устойчив на затухающем тестовом уравнении}.

\subsection{Решение неявного уравнения на каждом шаге}

В общем случае формула для $y_{i+1}$ не разрешается явно.
На каждом шаге нужно решить уравнение
\begin{equation*}
    G(v)=v-y_i-h f(x_{i+1},v)=0.
\end{equation*}
Простейший вариант --- \rconcept{простые итерации}:
\begin{equation*}
    y_{i+1}^{(k+1)}=y_i+h f(x_{i+1},y_{i+1}^{(k)}),
    \qquad k=0,1,\ldots.
\end{equation*}
Индекс $i$ нумерует узлы сетки, а $k$ --- внутренние итерации.
Начальным приближением может служить $y_i$ или явный прогноз Эйлера.

Если $f$ удовлетворяет условию Липшица по второму аргументу
с константой $L$, то достаточное условие сжатия этой итерации --- $hL<1$.
Для $f=-\lambda u$ ошибка внутренней итерации умножается на $-h\lambda$,
и требуется $h\lambda<1$.

\remark
\rnote{Устойчивость схемы и сходимость внутренней итерации --- разные
условия}. Неявная схема Эйлера устойчива на тестовом уравнении при любом
шаге, однако выбранный способ решения неявного уравнения может
потребовать ограничения $h$. Для больших шагов используют, например,
метод Ньютона или прямое решение линейной системы.

\section{Жёсткие системы дифференциальных уравнений}

\subsection{Разные масштабы затухания}

Пусть после выделения собственных мод линейная система содержит
уравнения
\begin{equation*}
    u_j'=-\lambda_j u_j,\qquad \lambda_j>0.
\end{equation*}
Каждая компонента затухает на своём характерном масштабе
\begin{equation*}
    \tau_j=\frac1{\lambda_j}.
\end{equation*}
Если $\max_j\lambda_j/\min_j\lambda_j$ велико, то в системе
одновременно присутствуют быстро и медленно затухающие моды.
Это характерный пример \rconcept{жёсткой системы}.

\begin{rbox}[breakable=false,left=18pt,right=18pt,top=14pt,bottom=10pt]{[]}
    \centering
    \begin{tikzpicture}[x=1.5cm,y=2.6cm,
        every node/.style={text=rtext,font=\small}]
        \draw[rtext,-{Stealth[length=2.2mm]}] (0,0)--(5.1,0) node[right] {$x$};
        \draw[rtext,-{Stealth[length=2.2mm]}] (0,0)--(0,1.2) node[left] {$u/u_0$};
        \draw[p3,line width=1.4pt,domain=0:4.7,samples=100]
            plot(\x,{exp(-0.55*\x)});
        \draw[red3,line width=1.4pt,domain=0:4.7,samples=100]
            plot(\x,{exp(-4*\x)});
        \node[p3,above right] at (1.8,0.372) {$e^{-\lambda_1 x}$};
        \node[red3,above right] at (0.32,0.278) {$e^{-\lambda_2 x}$};
        \node at (3.4,0.86) {$\lambda_2\gg\lambda_1$};
    \end{tikzpicture}
    \captionof{figure}{Быстрая мода почти исчезает раньше, чем существенно
    меняется медленная. Ограничение устойчивости явной схемы продолжает
    определяться быстрым масштабом.}
    \label{fig:ovf7_stiff_modes}
\end{rbox}

\subsection{Почему явная схема требует малого шага}

Для устойчивости явного метода Эйлера каждая мода должна удовлетворять
$h\lambda_j\le2$. Поэтому
\begin{equation*}
    h\le\frac2{\lambda_{\max}},\qquad
    \lambda_{\max}=\max_j\lambda_j.
\end{equation*}
Но интересующая динамика может происходить на медленном масштабе
$1/\lambda_{\min}$. Даже когда быстрая компонента стала очень малой,
слишком большой шаг будет усиливать возмущения этой компоненты,
в том числе ошибки округления.

Например, для $\lambda_1=1$ и $\lambda_2=1000$ явная схема требует
$h\le0.002$, хотя медленная компонента меняется на масштабе порядка
единицы. На отрезке длины $1$ приходится делать не менее $500$ шагов
уже по условию устойчивости.

Неявная схема позволяет выбирать шаг по требованиям точности
медленной динамики, а не только по быстрому масштабу затухания.
Это объясняет её пользу при решении жёстких задач.

\chapter{Численные методы линейной алгебры}
\rsubtitle{Метод Гаусса и прогонка для трёхдиагональных систем}

\section{Система линейных алгебраических уравнений}

Рассмотрим \rconcept{систему линейных алгебраических уравнений}
\begin{equation*}
    A\boldsymbol x=\boldsymbol b,\qquad \det A\ne0,
\end{equation*}
где
\begin{equation*}
    A=\begin{pmatrix}
        a_{11}&a_{12}&\cdots&a_{1n}\\
        a_{21}&a_{22}&\cdots&a_{2n}\\
        \vdots&\vdots&\ddots&\vdots\\
        a_{n1}&a_{n2}&\cdots&a_{nn}
    \end{pmatrix},\qquad
    \boldsymbol x=\begin{pmatrix}x_1\\x_2\\\vdots\\x_n\end{pmatrix},\qquad
    \boldsymbol b=\begin{pmatrix}b_1\\b_2\\\vdots\\b_n\end{pmatrix}.
\end{equation*}
Матрица $A$ и правая часть $\boldsymbol b$ заданы. Требуется найти
вектор неизвестных $\boldsymbol x$. Условие $\det A\ne0$
обеспечивает существование и единственность решения.

\section{Метод Гаусса}

\subsection{Идея исключения неизвестных}

\rconcept{Метод Гаусса} состоит из двух этапов.
На \rconcept{прямом ходе} последовательно исключают неизвестные
из нижних строк и приводят систему к верхнетреугольному виду:
\begin{equation*}
    \widetilde A\boldsymbol x=\widetilde{\boldsymbol b},\qquad
    \widetilde A=
    \begin{pmatrix}
        \widetilde a_{11}&\widetilde a_{12}&\cdots&\widetilde a_{1n}\\
        0&\widetilde a_{22}&\cdots&\widetilde a_{2n}\\
        \vdots&\ddots&\ddots&\vdots\\
        0&\cdots&0&\widetilde a_{nn}
    \end{pmatrix}.
\end{equation*}
На \rconcept{обратном ходе} неизвестные находят снизу вверх:
сначала $x_n$, затем $x_{n-1}$ и так далее до $x_1$.

\subsection{Прямой ход: один шаг алгоритма}

Пусть первые $k-1$ столбцов под диагональю уже обнулены.
Для каждой строки $i=k+1,\ldots,n$ вычислим множитель
\begin{equation*}
    m_{ik}=\frac{a_{ik}}{a_{kk}}.
\end{equation*}
Из строки $i$ вычитаем строку $k$, умноженную на $m_{ik}$:
\begin{align*}
    a_{ij}&\leftarrow a_{ij}-m_{ik}a_{kj},
        &&j=k+1,\ldots,n,\\
    b_i&\leftarrow b_i-m_{ik}b_k,\\
    a_{ik}&\leftarrow0.
\end{align*}
Здесь используются текущие, уже преобразованные коэффициенты.
Множитель $m_{ik}$ нужно сохранить \emph{до} обнуления $a_{ik}$.
Одинаковое преобразование применяется к матрице и правой части,
поэтому множество решений системы не меняется.

\subsection{Пример нескольких шагов}

Рассмотрим систему
\begin{equation*}
    \begin{cases}
        2x_1+x_2-x_3=1,\\
        4x_1+x_2+x_3=7,\\
        -2x_1+2x_2+x_3=4.
    \end{cases}
\end{equation*}
На первом шаге опорный элемент $a_{11}=2$.
Множители равны $m_{21}=2$ и $m_{31}=-1$:
\begin{equation*}
    \left(\begin{array}{rrr|r}
        2&1&-1&1\\
        4&1&1&7\\
        -2&2&1&4
    \end{array}\right)
    \longrightarrow
    \left(\begin{array}{rrr|r}
        2&1&-1&1\\
        0&-1&3&5\\
        0&3&0&5
    \end{array}\right).
\end{equation*}
На втором шаге $a_{22}=-1$, $m_{32}=-3$. Прибавляя к третьей
строке вторую, умноженную на $3$, получаем
\begin{equation*}
    \left(\begin{array}{rrr|r}
        2&1&-1&1\\
        0&-1&3&5\\
        0&0&9&20
    \end{array}\right).
\end{equation*}
Теперь система имеет треугольный вид.

\subsection{Обратный ход}

Последнее уравнение содержит только $x_n$:
\begin{equation*}
    x_n=\frac{\widetilde b_n}{\widetilde a_{nn}}.
\end{equation*}
После нахождения $x_{i+1},\ldots,x_n$ из строки $i$ получаем
\begin{equation*}
    x_i=\frac{\widetilde b_i-
        \displaystyle\sum_{j=i+1}^{n}\widetilde a_{ij}x_j}
        {\widetilde a_{ii}},\qquad i=n-1,\ldots,1.
\end{equation*}
В приведённом примере
\begin{equation*}
    x_3=\frac{20}{9},\qquad
    x_2=3x_3-5=\frac53,\qquad
    x_1=\frac{1-x_2+x_3}{2}=\frac79.
\end{equation*}

\subsection{Опорный элемент и стоимость}

Элемент $a_{kk}$, на который делят при исключении, называется
\rconcept{опорным}. Даже при $\det A\ne0$ он может оказаться нулевым
без перестановок строк. Например, у невырожденной матрицы
$\bigl(\begin{smallmatrix}0&1\\1&1\end{smallmatrix}\bigr)$ первый
диагональный элемент равен нулю.

При \rconcept{частичном выборе главного элемента} на шаге $k$
выбирают строку $p\ge k$, для которой $|a_{pk}|$ максимально,
и меняют местами строки $p$ и $k$ вместе с элементами правой части.
Это позволяет избежать деления на нуль и уменьшить влияние
малых опорных элементов на ошибки округления.

Для плотной матрицы прямой ход требует $O(n^3)$ операций,
обратный --- $O(n^2)$. Если матрица имеет специальную структуру,
её можно использовать для уменьшения стоимости решения.

\section{Метод прогонки для трёхдиагональной системы}

\subsection{Структура системы}

\rconcept{Трёхдиагональная матрица} имеет ненулевые элементы
только на главной диагонали и на двух соседних:
\begin{equation*}
    A=\begin{pmatrix}
        b_1&c_1&0&\cdots&0\\
        a_2&b_2&c_2&\ddots&\vdots\\
        0&a_3&b_3&\ddots&0\\
        \vdots&\ddots&\ddots&\ddots&c_{n-1}\\
        0&\cdots&0&a_n&b_n
    \end{pmatrix}.
\end{equation*}
Чтобы не путать главную диагональ $b_i$ с правой частью, последнюю
здесь обозначим через $d_i$. Система имеет вид
\begin{equation*}
    a_i x_{i-1}+b_i x_i+c_i x_{i+1}=d_i,
    \qquad i=1,\ldots,n,
\end{equation*}
где $a_1=0$ и $c_n=0$, то есть отсутствующие крайние слагаемые опускаются.

\rconcept{Метод прогонки}, или алгоритм Томаса, представляет собой
метод Гаусса с учётом этой структуры. При исключении достаточно
изменять одну диагональ и правую часть.

\subsection{Прямой ход прогонки}

Начнём с первой строки:
\begin{equation*}
    \widetilde b_1=b_1,\qquad \widetilde d_1=d_1.
\end{equation*}
Чтобы исключить $x_{i-1}$ из строки $i$, используем уже
преобразованную строку $i-1$:
\begin{equation*}
    \widetilde b_{i-1}x_{i-1}+c_{i-1}x_i=\widetilde d_{i-1}.
\end{equation*}
Множитель исключения равен
\begin{equation*}
    m_i=\frac{a_i}{\widetilde b_{i-1}}.
\end{equation*}
Вычитая предыдущую строку, умноженную на $m_i$, получаем
\begin{align*}
    \widetilde b_i&=b_i-\frac{a_i c_{i-1}}{\widetilde b_{i-1}},\\
    \widetilde d_i&=d_i-\frac{a_i\widetilde d_{i-1}}
        {\widetilde b_{i-1}},\qquad i=2,\ldots,n.
\end{align*}
Поддиагональный элемент обнуляется, а верхняя диагональ
$c_i$ остаётся без изменения. В результате система становится
двухдиагональной:
\begin{equation*}
    \widetilde b_i x_i+c_i x_{i+1}=\widetilde d_i.
\end{equation*}

\subsection{Обратный ход прогонки}

В последней строке нет $x_{n+1}$, поэтому
\begin{equation*}
    x_n=\frac{\widetilde d_n}{\widetilde b_n}.
\end{equation*}
Далее движемся справа налево:
\begin{equation*}
    x_i=\frac{\widetilde d_i-c_i x_{i+1}}{\widetilde b_i},
    \qquad i=n-1,n-2,\ldots,1.
\end{equation*}
В отличие от общего обратного хода Гаусса, здесь для нахождения
$x_i$ используется только одно уже найденное значение $x_{i+1}$.
На каждой строке выполняется фиксированное число действий,
поэтому прогонка требует \rresult{$O(n)$ операций} и $O(n)$ памяти.

\remark
Формулы требуют ненулевых знаменателей $\widetilde b_i$.
Одной невырожденности исходной матрицы для прогонки без перестановок
недостаточно. Достаточное условие отсутствия нулевых опорных элементов ---
строгое диагональное преобладание
\begin{equation*}
    |b_i|>|a_i|+|c_i|\qquad\text{для всех }i.
\end{equation*}
Это достаточное, а не необходимое условие применимости алгоритма.

\subsection{Откуда возникают трёхдиагональные матрицы}

Такие системы появляются при разностной аппроксимации второй
производной. Например, уравнение диффузии
\begin{equation*}
    \frac{\partial u}{\partial t}
    =D\frac{\partial^2u}{\partial x^2},\qquad D>0,
\end{equation*}
в стационарном случае сводится к $u''(x)=0$.
На равномерной сетке вторая производная заменяется формулой
\begin{equation*}
    \frac{y_{i+1}-2y_i+y_{i-1}}{h^2}=0.
\end{equation*}
Каждое уравнение связывает только три соседних значения.

\begin{rbox}[breakable=false,left=18pt,right=18pt,top=14pt,bottom=10pt]{[]}
    \centering
    \begin{tikzpicture}[x=2.5cm,y=1cm,
        every node/.style={text=rtext,font=\small}]
        \draw[rtext,-{Stealth[length=2.2mm]}] (-0.25,0)--(3.25,0) node[right] {$x$};
        \foreach \xx/\lab in {0/{x_{i-1}},1.5/{x_i},3/{x_{i+1}}}{
            \fill[p3] (\xx,0) circle(2pt);
            \node[below=5pt] at (\xx,0) {$\lab$};
        }
        \node[above=7pt] at (0,0) {$y_{i-1}$};
        \node[above=7pt] at (1.5,0) {$y_i$};
        \node[above=7pt] at (3,0) {$y_{i+1}$};
        \draw[rtext!65,{Stealth[length=1.5mm]}-{Stealth[length=1.5mm]}]
            (0,-0.85)--(1.5,-0.85) node[midway,fill=rbg] {$h$};
        \draw[rtext!65,{Stealth[length=1.5mm]}-{Stealth[length=1.5mm]}]
            (1.5,-0.85)--(3,-0.85) node[midway,fill=rbg] {$h$};
    \end{tikzpicture}
    \captionof{figure}{Трёхточечный шаблон второй производной:
    строка матрицы содержит коэффициенты при $y_{i-1}$, $y_i$, $y_{i+1}$.}
    \label{fig:ovf7_three_point_stencil}
\end{rbox}

Пусть заданы граничные значения $y_0=U_L$ и $y_{n+1}=U_R$,
а неизвестны $y_1,\ldots,y_n$. После умножения уравнений на $-h^2$
и переноса граничных значений вправо получаем
\begin{equation*}
    \begin{pmatrix}
        2&-1&0&\cdots&0\\
        -1&2&-1&\ddots&\vdots\\
        0&\ddots&\ddots&\ddots&0\\
        \vdots&\ddots&-1&2&-1\\
        0&\cdots&0&-1&2
    \end{pmatrix}
    \begin{pmatrix}y_1\\y_2\\\vdots\\y_{n-1}\\y_n\end{pmatrix}
    =\begin{pmatrix}U_L\\0\\\vdots\\0\\U_R\end{pmatrix}.
\end{equation*}
Во внутренних строках правая часть нулевая, но в крайних строках
в неё входят известные граничные значения.

Для этой матрицы $\widetilde b_1=2$ и
\begin{equation*}
    \widetilde b_i=2-\frac1{\widetilde b_{i-1}}
    =\frac{i+1}{i}>0.
\end{equation*}
Поэтому прогонка работает без нулевых знаменателей, хотя во внутренних
строках диагональное преобладание нестрогое.

\end{document}
